\chapter{Title bars and notifications}
\label{ch:titles}

\begin{aside}
A long-running TUI deserves a name in the tab bar and the
ability to nudge the user when something interesting happens.
Both are one escape away.
\end{aside}

\section{Setting the title}

\code{ESC ] 0 ; TITLE BEL} sets both the window title and the
icon title. \code{ESC ] 1 ; TITLE BEL} sets just the icon
title; \code{ESC ] 2 ; TITLE BEL} sets just the window title.

\maya{}'s \code{set\_title} call uses OSC 0 (both) by default
because most apps want the same string in both places.

\section{Dynamic titles}

A common pattern: the title reflects the current document or
state.

\begin{lstlisting}
ctx.effect([&] {
    auto path = ctx.get(current_file);
    auto dirty = ctx.get(is_dirty);
    maya::set_title(
        path + (dirty ? " *" : "")
    );
});
\end{lstlisting}

The effect re-runs when either signal changes; the title is
always in sync.

\section{Title and tab order}

In tmux, the pane title is set via the same escape but
displayed in the status bar (if configured). In a tabbed
terminal, the title is the tab label. Setting it changes how
the user identifies your app among many.

\section{OSC 7: current working directory}

\code{ESC ] 7 ; file://hostname/path BEL} tells the terminal
the current working directory. Opening a new tab inherits it.
\maya{} emits OSC 7 on startup and whenever the app's notion
of cwd changes (file pickers, for example).

\section{OSC 9: notification}

\code{ESC ] 9 ; message BEL} (iTerm2 extension) raises an OS
notification. The user gets a native popup with your message.
Supported by iTerm2, Wezterm, Ghostty.

\maya{} provides \code{notify(message)} that emits OSC 9. For
terminals that don't support it, the escape is silently
ignored.

\section{Bell}

\code{BEL} (\code{\textbackslash a}, byte \code{0x07}) rings
the terminal bell. Most terminals show a visual or audible
indicator. \maya{} exposes \code{bell()} that emits one byte.

Bells are obnoxious if overused. Reserve for ``something
needs attention now'' moments, not for every notification.

\begin{funfact}
The bell predates terminals: it was a real bell on a
teletype, triggered by the same byte. ASCII \code{0x07} is
literally ``ring the bell''.
\end{funfact}

\section{Activity indicators}

In tmux, the bell triggers an activity flag on the pane. The
user sees ``activity'' in their status bar even when looking
at another pane. Use bells (or OSC 9 notifications) to make
background work visible.

\section{Window icon}

OSC 1 sets the icon title, which becomes the icon-bar label
when the window is minimised on traditional X11 setups. Not
many users see this anymore, but the escape is still there.

\section{Window manipulation}

\code{ESC [ 8 ; H ; W t} resizes the terminal (if the
terminal allows). \code{ESC [ 14 t} reports pixel size.
\code{ESC [ 18 t} reports cell size. \code{ESC [ 22 ; 0 t}
saves the title to a stack; \code{ESC [ 23 ; 0 t} pops.

\maya{} pushes the title on startup and pops on exit, so the
user's previous title is restored. If a child process (e.g.
a shell spawned via PTY in your app) changes the title, the
pop brings it back.

\section{Coloured tabs}

\code{ESC ] 6 ; 1 ; bg ; red ; brightness ; 255 BEL} sets the
tab colour in iTerm2. Other terminals have their own
extensions. \maya{} does not abstract these; if you want to
target a specific terminal, emit its sequence directly.

\section{Progress in the title}

A useful pattern for batch jobs:

\begin{lstlisting}
ctx.effect([&] {
    int pct = ctx.get(progress);
    maya::set_title(
        "myapp - " + std::to_string(pct) + "%"
    );
});
\end{lstlisting}

The user sees the progress in their tab bar even when the
app is not focused.

\section{Recap}

\begin{recap}
\begin{itemize}[leftmargin=*]
  \item OSC 0, 1, 2 set the title.
  \item OSC 7 announces the cwd.
  \item OSC 9 raises a notification (iTerm2 extension).
  \item Bell is a one-byte attention prod.
  \item \maya{} pushes/pops the title around its lifetime.
\end{itemize}
\end{recap}

\section*{Workshop}

\begin{enumerate}[leftmargin=*]
  \item Bind the title to your app's current state (file
    name, view, percent complete).
  \item Add OSC 9 notifications for long-running operations
    when they complete.
  \item Push/pop the title from a sub-mode (e.g., entering
    a settings dialog).
\end{enumerate}
